// Const Values — demonstrate const declarations for color grading
//
// The const keyword declares a compile-time constant. Once declared,
// the value cannot be reassigned — attempting to write to a const
// variable will produce an error. This is useful for defining fixed
// reference values that should never change during evaluation.
//
// This example performs a 3-way color grade: shadows, midtones, and
// highlights each receive a different color tint.

f$shadow_thresh = 0.3;    // luminance below this is shadows
f$high_thresh = 0.7;      // luminance above this is highlights

// ── Const color definitions (cannot be reassigned) ──────────────
const vec3 shadow_tint = vec3(0.15, 0.18, 0.35);    // cool blue shadows
const vec3 mid_tint    = vec3(0.30, 0.28, 0.20);     // warm amber mids
const vec3 high_tint   = vec3(1.0, 0.95, 0.85);      // warm white highlights
const float tint_strength = 0.4;                      // how much tint to apply

// NOTE: the following would cause an error because const prevents reassignment:
//   shadow_tint = vec3(1.0, 0.0, 0.0);   // ERROR — cannot reassign const

vec3 src = @image;
float lum = luma(src);

// Determine which zone this pixel falls into
float s_weight = 1.0 - smoothstep($shadow_thresh - 0.05, $shadow_thresh + 0.05, lum);
float h_weight = smoothstep($high_thresh - 0.05, $high_thresh + 0.05, lum);
float m_weight = 1.0 - s_weight - h_weight;

// Blend the zone tints
vec3 tint = shadow_tint * s_weight + mid_tint * m_weight + high_tint * h_weight;

// Apply: lerp between the source and the tinted version
@OUT = lerp(src, src * tint / max(luma(tint), 0.01), tint_strength);
